\scanpage{511}
\mspage{261}

-- Here one uses the fact that the object $X(\mathbb C)\setminus\{0,1,\infty\}$ of the category of \unsure{covering\textsuperscript{s}} with maximum ramification \unsure{imposed|imposed}, attains precisely a maximum, at the points $1,\infty$ where the condition comes in.
One finds that the corresponding functor
\[
\tag{7 bis}
\Pi'_{03}\text{-}\mathrm{Ens}\longrightarrow(\Pi_{0,3}\text{-}\mathrm{Ens})
\]
has the exactness properties of a functor associated to a morphism of Galois topoi
\diagram{\begin{tikzcd}[ampersand replacement=\&]
\overset{\inserted{\unsure{B}}}{\unsure{B|D}}_{\Pi_{03}} \arrow[r] \& B_{\Pi'_{0,3}}
\end{tikzcd}}
(\unsure{commutes with} $\underset{\longrightarrow}{\lim}$ \unsure{of any kind}, left exact) and is therefore defined by a $(\Pi_{03}^{\deleted{\unsure{\wedge}}},\Pi'_{03})$ set which is a (right) torsor for $\Pi'_{03}$. The choice of an origin for this torsor then also defines a corresponding homomorphism $\Pi_{0,3}\longrightarrow\Pi'_{03}$.

\margin{This bitorsor is also the set of ($Y(\mathbb C)\setminus$\deleted{\illegible{2}}\inserted{\unsure{$\{0\}$}}) \unsure{isomorphisms|morphisms} of the ``universal covering with \unsure{ramification} imposed'' of this space, onto the \deleted{\unsure{with ramification}} \deleted{\unsure{ordinary}} ordinary universal covering of $X(\mathbb C)\setminus\{0,1,\infty\}$ \ldots}

Let us return to the arithmetic situation over $\mathbb Q$, where one has not only the geometric profinite fundamental groups $\widehat{\Pi}_{03}$, $\widehat{\Pi}'_{03}$, but also extensions
\[
\tag{8}
\left\{\begin{aligned}
1&\unsure{\longrightarrow}\widehat{\Pi}_{03}\longrightarrow\mathcal E_{0,3}^{\deleted{\unsure{\wedge}}}\longrightarrow\Gamma_{\deleted{\unsure{\mathbb Q}}}\longrightarrow1\\
1&\unsure{\longrightarrow}\widehat{\Pi}'_{03}\unsure{\longrightarrow}\mathcal E_{0,3}^{\prime\deleted{\unsure{\wedge}}}\unsure{\longrightarrow}\Gamma\unsure{\longrightarrow}1
\end{aligned}\right.,
\]
interpreted as fundamental groups of $\unsure{U}_{03\deleted{\mathbb Q}}$, resp. of $\unsure{U}_{03}$ with imposed ramification $2\{1\}+3\{\infty\}$. The arguments

\scanpage{512}

above extend to this setting, and thus provide a homomorphism \deleted{outer} \inserted{of extensions}
\diagram{\begin{tikzcd}[ampersand replacement=\&]
1 \arrow[r] \& \widehat{\Pi}_{03} \arrow[r] \arrow[d] \& \mathcal E_{03} \arrow[r] \arrow[d] \& \Gamma \arrow[r] \arrow[d,equal] \& 1 \\
1 \arrow[r] \& \widehat{\Pi}'_{03} \arrow[r] \& \mathcal E'_{03} \arrow[r] \& \Gamma \arrow[r] \& 1
\end{tikzcd}\quad(9)},
modulo conjugation by an inner automorphism of $\Pi'_{03}$.

\margin{(9 bis) $\widehat{\Pi}_{0,3}\longrightarrow\widehat{\Pi}'_{0,3}$ fits into an exact sequence
\[
\boxed{1\longrightarrow\widehat{\Pi}_{0,3}\longrightarrow\widehat{\Pi}'_{0,3}\longrightarrow\unsure{\mathfrak S}_3\longrightarrow1}
\]
and ditto for the discrete groups
\[
\boxed{1\longrightarrow\Pi_{0,3}\longrightarrow\Pi'_{0,3}\longrightarrow\unsure{\mathfrak S}_3\longrightarrow1}.
\]}

That said, let $K$ be the kernel of the outer operation of $\Gamma$ on $\Pi'_{03}$. One then sees that \deleted{$K$ operates} the outer operation of $K$ on the open subgroup $\widehat{\Pi}_{03}$ takes place through a finite group (in fact, through \inserted{the intermediary of} $\widehat{\Pi}'_{03}/\widehat{\Pi}_{03}\simeq\unsure{\mathfrak S}_3$ \ldots) -- which implies, by theorem 1 of the \S\ above, that the group $K$ itself is finite. Thus by Artin one has $K=1$ or $K=(1,\tau)$, $\tau$ complex conjugation. But since $\chi(\tau)=-1$ it is obvious that the outer operation of $\tau$ on $\Pi'_{03}$ \inserted{(nor even on $(\Pi'_{03})_{\unsure{\mathrm{ab}}}^{*}$ \deleted{\illegible{2}})} is not trivial\ldots
\margin{$*\simeq\unsure{\mathbb Z/2\mathbb Z+\mathbb Z/3\mathbb Z}$}
\deleted{\unsure{Which}} This proves the proposition.

\scanpage{513}
\mspage{262}

I now feel like interpreting $\Pi_{03}^{\unsure{\prime}}$ as $\mathcal T^+_{1,1}/\deleted{\text{subgroup}}\inserted{\text{center}}$ (isomorphic \unsure{to} $\{\pm1\}$\unsure{),|,} and ditto for $\widehat{\Pi}'_{03}\simeq\widehat{\mathcal T}_{11}/\{\pm1\}$, \unsure{isomorphism|isomorphisms} which is compatible with the outer action of $\Gamma$, One has (for the record)
\[
\tag{10}
\mathcal T_{1,1}\xrightarrow{\sim}\mathrm{Gl}(2,\mathbb Z),
\]
the isomorphism being obtained as follows:
\[
\tag{11}
\mathcal T_{1,1}\xrightarrow[\substack{\text{plugging}\\\text{the hole}}]{\mathrm{iso}}\mathcal T_{1,0}\xrightarrow{\mathrm{iso}}\mathrm{Autext}\left(\underbrace{\Pi_{10}}_{\substack{\wr\\\mathbb Z^2}}\right)\simeq\mathrm{Gl}(2,\mathbb Z)
\]

One thus has
\[
\tag{12}
\mathcal T^+_{11}\simeq\mathrm{Sl}(2,\mathbb Z)
\]
and
\[
\tag{13}
\mathrm{Cent}(\mathcal T_{11})\simeq(\pm1)
\]
(identifies with the subgroup of homotheties of $\mathrm{Sl}(2,\mathbb Z)$). It is known that one has \unsure{in} $\mathrm{Sl}(2,\mathbb Z)/(\pm1)$ two generators $\unsure{\lambda_3,\lambda_2}$, satisfying respectively
\[
\tag{14}
\lambda_3^3=1,\quad\lambda_2^2=1
\]
and such that these relations are a \uline{presentation} of $\mathrm{Sl}(2,\mathbb Z)/(\pm1)$. Unless I am mistaken, these \unsure{elements|elements} come from \unsure{elements|elements} $u_6,u_4$ of $\mathrm{Sl}(2,\mathbb Z)$ itself, satisfying
\[
\tag{15}
u_6^6=1,\quad u_4^4=1
\]

\scanpage{514}

and which correspond to the only elements of order 6 and of order 4 (up to conjugation and taking inverses). In fact, \unsure{every} element of finite order of $\mathrm{Sl}(2,\mathbb Z)$ is conjugate to a power \unsure{of|of} \deleted{a} $u_4$, or $u_6$. The \unsure{subgroups} \unsure{generated} resp. by $u_4$ and $u_6$, cyclic of order 4 and 6, are the groups of \unsure{automorphisms} of the two exceptional elliptic curves (called ``\unsure{anharmonic}'') (in characteristic 0).

\deleted{For} I have no formula at hand for ``placing'' \deleted{\unsure{$u_4$}} $u_4$ and $u_6$, so that they generate the group $\mathrm{Sl}(2,\mathbb Z)$\deleted{\unsure{$/\pm1$}} \unsure{modulo} its center (which probably implies that they generate $\mathrm{Sl}(2,\mathbb Z)$ --- but \unsure{one will pay attention|well. Pay attention} to the additional relation. \deleted{being added \unsure{to} \inserted{(15)} \unsure{of the} fundamental relations}
\[
(16)\qquad u_4^2=u_6^3\qquad\unsure{\varepsilon}
\]
(this is \deleted{\illegible{2}} \unsure{precisely} the element $-1$ of the center of $\mathrm{Sl}(2,\mathbb Z)$). Perhaps I do not need it (the explicit formula for $u_4,u_6$). One then \unsure{will define|defines}

\scanpage{515}
\mspage{263}

\[
(17)\qquad \Pi'_{03}\xrightarrow{\sim}\mathrm{Sl}(2,\mathbb Z)/(\pm1)
\]
by
\[
(18)\qquad
\begin{array}{rcl}
\ell'_1&\longmapsto&\lambda_2\\
\ell'_\infty&\longmapsto&\lambda_3
\end{array}
\]
\[
\bigl[\text{thus}\qquad\ell'_0\longmapsto(\lambda_3\lambda_2)^{-1}
=\lambda_2\lambda_3^{-1}=\lambda_2\lambda_3^2\bigr],
\]
\unsure{where} $\ell'_0,\ell'_1,\ell'_\infty$ are the generators of $\Pi'_{03}$ corresponding to those $\ell_0,\ell_1,\ell_\infty$ of $\Pi_{03}$, with the defining relations
\[
(19)\qquad\ell'_\infty\ell'_1\ell'_0=1\qquad
\ell_1^{\prime 2}=1,\quad\ell_\infty^{\prime 3}=1
\]
so that $\Pi'_{03}$ is indeed the group \deleted{free} with generators $\ell'_1,\ell'_\infty$ satisfying: $\ell_1^{\prime 2}=1,\ell_\infty^{\prime 3}=1$.

I now want to convince myself that the homomorphism corresponding to (17)
\[
(20)\qquad\widehat{\Pi}'_{03}\longrightarrow
\mathrm{Sl}(2,\mathbb Z)^\wedge/\pm1\simeq
\widehat{\mathcal T}_{1,1}^{+}/(\pm1)
\]
is compatible with the outer operation of $\Gamma=\Pi_{\mathbb Q}$. For this, I feel like defining directly the composite \unsure{with} $\unsure{\widehat{\Pi}_{0,3}|\Pi_{0,3}}\longrightarrow\widehat{\Pi}'_{03}$
\[
\text{\deleted{\unsure{(21)}}}\qquad\widehat{\Pi}_{03}\longrightarrow\widehat{\mathcal T}_{1,1}^{+}/\pm1
\]
and even \unsure{lifting} it to a homomorphism
\[
(21)\qquad\widehat{\Pi}_{03}\longrightarrow\widehat{\mathcal T}_{11}^{+}\;.
\]

\scanpage{516}

At the level of discrete groups, one does define
\[
(22)\qquad\Pi_{03}\longrightarrow\mathrm{Sl}(2,\mathbb Z)
\]
\[
\left\{
\begin{array}{rcl}
\ell_1&\longmapsto&u_4\\
\ell_\infty&\longmapsto&u_6\\
\ell_0&\longmapsto&(u_6u_4)^{-1}=u_4^{-1}u_6^{-1}
\end{array}
\right.
\]
and $\mathrm{Sl}(2,\mathbb Z)$\inserted{$\simeq\mathcal T_{11}^{+}$} identifies with the quotient of $\Pi_{0,3}$ by the invariant \unsure{subgroup} \unsure{generated} by the elements
\[
\ell_1^4,\ \ell_\infty^6,\ (\ell_1^2)(\ell_\infty^3)^{-1}\;.
\]

Besides, $\Pi_{03}$ can be interpreted as a Teichmüller group $\mathcal T_{0,4}^{!+}$, and (22) as a homomorphism
\margin{indeed, $\mathcal T_{g,\nu}^{!}$ is an extension of $\mathcal T_{g,\nu-1}^{!}$ by $\Pi_{g,\nu-1}$ (if $(g,\nu-1)$ \unsure{anabelian}) and ditto for the $\mathcal T^{+}$, $\mathcal T_{0,3}^{!+}=(1)$!}
\[
\text{\deleted{\unsure{(24)}}}\ (23)\qquad\mathcal T_{0,4}^{!+}\longrightarrow\mathcal T_{11}^{+}
\]
which I suspect extends to a homomorphism
\[
(24)\qquad
\diagram{\begin{tikzcd}[ampersand replacement=\&,column sep=large]
\mathcal T_{0,4}^{!}\arrow[r]\arrow[d,phantom,"\wr" description]\&\mathcal T_{11}\arrow[d,phantom,"\wr" description]\\
\begin{gathered}
\text{extension of }\{1,\unsure{\tau}\}\text{ by}\\
\Pi_{0,3},\ \text{which is none}\\
\text{other than the}\\
\text{cartographic group}\\
\text{triangulated weighted}\\
\text{unoriented}
\end{gathered}
\&\mathrm{Gl}(2,\mathbb Z)
\end{tikzcd}}
\]
\margin{not yet defined}

I interpret the first \unsure{side} \inserted{of (23)} as the $\Pi_1$ ``transcendental geometric'' \unsure{of the|of the} \unsure{modular|modular} topos

\scanpage{517}
\mspage{264}

$\mathbb{M}_{0,4}^{!}$, classifying \deleted{rational curves} \inserted{projective lines} with four distinct points numbered from $1$ to $4$; and the second \unsure{side|right-hand side} of (23) is the transcendental $\pi_1$ of the modular topos $M_{1,1}$, classifying elliptic curves (with a fixed origin). One should therefore be able to define, at the level of \deleted{schemes} modular multiplicities over $\mathbb{Q}$
\[
(25)\qquad (\mathbb{M}_{0,4}^{!})_{\mathbb{Q}} \longrightarrow
(\mathbb{M}_{1,1}^{\deleted{\unsure{!}}})_{\mathbb{Q}}
\]
which would give rise to (23). \deleted{This \unsure{also means}}

But one sees at once that one has a canonical isomorphism
\[
(26)\qquad \mathbb{M}_{0,4}^{!}\simeq U_{0,3}
\]
(I drop the subscripts $\mathbb{Q}$), and giving (25) thus also amounts to giving a family of elliptic curves over $U_{0,3}$.

But \unsure{if} $\lambda$ is the \unsure{``variable'' $\in U_{0,3}$''}, taking ``the'' quadratic covering $E_\lambda$ of \deleted{\illegible{1}} $\mathbb{P}^1$ ramified at $0,1,\infty,\lambda$,

\scanpage{518}

one finds an elliptic curve, with four marked points \deleted{resp.} above $0,1,\infty,\lambda$ --- of which one can take the point above $0$ as origin --- then the other points are the three points of order $2$, indexed resp. by $1,\infty,\lambda$.

More precisely: over any base scheme \inserted{$S$}, except that $2$ should be invertible there (residual characteristics \unsure{all} $\ne 2$), \deleted{the} \inserted{one has a} functor which goes \deleted{from the} \inserted{from the category of systems $(E,\alpha,\beta\deleted{,\unsure{\gamma}})$ of a} \unsure{relative|relative} \unsure{elliptic|elliptic} \unsure{curves|curve} (\unsure{homogeneous|homogeneous}) \deleted{\unsure{with} \illegible{1}} over $S$, and \deleted{\unsure{whose}} $\alpha,\beta\deleted{,\unsure{\gamma}}$ \deleted{three} \inserted{two} sections of $E$ which form a basis of ${}_2E$ \unsure{as} a \inserted{locally constant} scheme in $\mathbb{F}_2$-Modules (a local system of $\mathbb{F}_2$-modules) on $S$ [\unsure{giving} rise to: $\gamma=\alpha+\beta$, as a third rogue --- so that $(\beta,\gamma)$\unsure{,|.} \deleted{and \unsure{$(\gamma,\alpha)$}} \inserted{etc} are in fact also bases) --- to the category of projective-line bundles $P$ over $S$, with four sections \inserted{$u_0,u_1,u_\infty,\lambda$} \unsure{marked} distinct at every point, by associating to $E$ the quotient $E/\pm1$, \unsure{equipped} with the sections $u_0,u_1,u_\infty,\lambda$ which are the images resp. of $0,\alpha,\beta,\gamma=\alpha+\beta$

\margin{it amounts to the same thing to say that \unsure{one has|the} three sections $\alpha,\beta,\gamma$ of ${}_2E$ over $S$, which \unsure{are} $\forall s\in S$, \deleted{\unsure{everywhere}} distinct and \unsure{hence|everywhere} $\ne1$.}

\scanpage{519}
\mspage{265}

\margin{The pages which follow are needlessly complicated, with the introduction of the $\unsure{M|\mathbb M}'_{1,1}$, $\mathcal{T}'_{1,1}$ \unsure{etc.} -- it is enough to go at it \unsure{brutally|simply} with the elliptic curve $y^2=\sqrt{x(x-1)(x-\lambda)}$, over $U_{0,3}$, \unsure{to} \illegible{1} \unsure{have} $U_{0,3}\longrightarrow \unsure{M|\mathbb M}_{1,1}$ cf. below\ldots}

and this functor is \uline{almost} an equivalence of categories. What it lacks in order to be one, is that for a relative elliptic curve $E$, the symmetry \inserted{$x\mapsto -x$} of $E$ -- which is a nontrivial automorphism -- operates \unsure{trivially} on the corresponding object $E/\pm1$. So one \unsure{must} take \unsure{the|the} \deleted{fibered category} \inserted{stack} over $(\mathrm{Sch})$, obtained from that of elliptic curves by starting by taking $\unsure{\mathrm{Isom}}(E,F)'=\unsure{\mathrm{Isom}}(E,F)/\pm1$, then one will take the stack associated to a prestack (the objects over $S$ are the ``relative elliptic curves up to symmetry over $S$''). \unsure{This|This|The} stack is representable by a modular multiplicity $\unsure{M|\mathbb M}'_{1,1}$, \deleted{\unsure{obtained from} \illegible{1}} having as its geometric fundamental group $\mathrm{Sl}(2,\mathbb{Z})^{\wedge}/\pm1$ precisely, obtained \deleted{from the} for example from the modular varieties with Jacobi rigidification of level $n$, $\unsure{M|\mathbb M}_{11}[n]$ -- with group $\mathrm{Sl}(2,\mathbb{Z}/n\mathbb{Z})$ operating on it, so that
\[
 (27)\qquad \unsure{M|\mathbb M}_{11}\simeq \bigl(\unsure{M|\mathbb M}_{1,1}[n],\mathrm{Sl}(2,\mathbb{Z}/n\mathbb{Z})\bigr)
\]
\unsure{moreover|one notes} that \inserted{the central subgroup $\unsure{(\pm1)|\{\pm1\}}$ of} $\mathrm{Sl}(2,\mathbb{Z}/n\mathbb{Z})$ operates \unsure{trivially} on $\unsure{M|\mathbb M}_{11}[n]$, so one can let the quotient group $\mathrm{Sl}(2,\mathbb{Z}/n\mathbb{Z})'\,\deleted{\illegible{1}}=\mathrm{Sl}(2,\mathbb{Z}/n\mathbb{Z})/\pm1$ operate, and set

\scanpage{520}

\[
 (28)\qquad \unsure{M|\mathbb M}'_{11}=\bigl(\unsure{M|\mathbb M}_{11}[n],\mathrm{Sl}(2,\mathbb{Z}/n\mathbb{Z})'\bigr)
\]
\unsure{(which|Which} \unsure{manifestly$^{\mathrm t}$|manifestly} does not depend on the choice of $n$, $n\geq2$). One thus finds a morphism
\[
 (29)\qquad
 \diagram{\begin{tikzcd}[ampersand replacement=\&,row sep=small]
 \unsure{M|\mathbb M}^!_{0,4} \arrow[r] \arrow[d,equal] \& \unsure{M|\mathbb M}'_{11} \\
 U_{0,3} \&
 \end{tikzcd}}
\]
which makes $\unsure{M|\mathbb M}^!_{0,4}$ a Galois covering with group $\unsure{\mathfrak{S}_3}$ of $\unsure{M|\mathbb M}'_{11}$, and more precisely
\[
 (30)\qquad
 \diagram{\begin{tikzcd}[ampersand replacement=\&,column sep=large]
 \unsure{M|\mathbb M}^!_{0,4} \arrow[r,dash,"\sim"] \&
 \unsure{M|\mathbb M}_{1,1}^{\prime\,\deleted{\illegible{1}}}[2]
 \arrow[r,"{\text{Galois covering with group }\mathrm{Sl}(2,\mathbb{F}_2)\simeq\unsure{\mathfrak{S}_3}}"'] \&
 \unsure{M|\mathbb M}'_{11}\,\deleted{\illegible{1}}
 \end{tikzcd}}
\]
the geometric fundamental group of $\unsure{M|\mathbb M}_{11}[2]$ \unsure{being} moreover isomorphic to the kernel of \inserted{the homomorphism} $\mathrm{Sl}(2,\mathbb{Z})^{\prime\wedge}\longrightarrow\mathrm{Sl}(2,\mathbb{Z}/2\mathbb{Z})$ which factors the canonical homomorphism $\mathrm{Sl}(2,\mathbb{Z})^{\wedge}_{\deleted{\illegible{1}}}\longrightarrow\mathrm{Sl}(2,\mathbb{Z}/2\mathbb{Z})$.

Passing to the fundamental groups for $\unsure{M|\mathbb M}^!_{0,4}=U_{0,3}$ and $\unsure{M|\mathbb M}'_{11}$, one finds a \unsure{homomorphism|homomorphism} of exact sequences

\deleted{\unsure{(31)}}
\[
 (31)\qquad
 \diagram{\begin{tikzcd}[ampersand replacement=\&,column sep=small,row sep=small]
 \& \inserted{\unsure{\mathcal{T}_{0,4}^{!+}|\widehat{\mathcal{T}}_{0,4}^{!+}}}
   \arrow[d,phantom,"\inserted{\Vert}" description]
 \& \inserted{\unsure{N_{0,4}|\mathcal{N}_{0,4}}}
   \arrow[d,phantom,"\inserted{\Vert}" description] \& \& \\
 1 \arrow[r] \& \widehat{\Pi}_{0,3} \arrow[r] \arrow[d]
 \& \mathcal{E}_{0,3} \arrow[r] \arrow[d]
 \& \Gamma \arrow[r,dash] \arrow[d,"\Vert"] \& 1 \\
 1 \arrow[r] \& \widehat{\mathcal{T}}_{1,1}^{\prime+} \arrow[r]
   \arrow[d,phantom,"\inserted{\wr}" description]
 \& \mathcal{E}'_{11} \arrow[r] \& \Gamma \arrow[r] \& 1 \\
 \& \inserted{\widehat{\mathcal{T}}_{11}^{+}/\unsure{\{\pm1\}|(\pm1)}
   \simeq\mathrm{Sl}(2,\mathbb{Z})^{\wedge}/\pm1} \& \& \&
 \end{tikzcd}}
\]
